\section{Project Overview}
\label{section:intro}

\subsection{Profile}

\paragraph{Track: 2. Dataset and Benchmark / 3. Application}
An evaluation study: can \emph{existing, publicly available} faithfulness-verification
techniques be adopted as a verification layer for a retrieval-augmented QA system?

\paragraph{Abstract and project goal:}
Retrieval-augmented generation (RAG) grounds large language model (LLM) answers in
retrieved documents, but a grounded \emph{citation} does not guarantee a grounded
\emph{claim}: an answer can cite a source and still assert something the cited passage
does not support. A growing body of \emph{published} work offers ways to measure this
``faithfulness''---LLM-as-judge scoring (RAGAS~\citep{es2024ragas}), sampling consistency
(SelfCheckGPT~\citep{manakul2023selfcheckgpt}), atomic-fact and claim-level methods
(FActScore~\citep{min2023factscore}, VeriScore~\citep{song2024veriscore}),
NLI-based factual consistency (TRUE~\citep{honovich2022true}), and purpose-built detectors
benchmarked in~\citet{sardana2025realtime}. This project \textbf{evaluates the adoption of
these published techniques} as a post-generation verification layer on a RAG QA system:
we segment answers into claims, check each claim against its cited source using the
published scorers above, and benchmark their agreement with human labels, cost, and
reproducibility. The goal is a practical recommendation---\emph{which published
verification technique(s) are worth adopting, and under what trade-offs}---evaluated on
public benchmarks.

\paragraph{TL;DR (``Too Long; Didn't Read"):}
An evaluation of existing published faithfulness-verification techniques, applied and
benchmarked as a claim-checking layer on a RAG QA system.

\subsection{Motivation}
\begin{itemize}
    \item \textbf{Why is the problem interesting?}
    RAG is the default recipe for grounding LLMs, yet deployed systems still emit
    ``plausible but unsupported'' statements. Many verification techniques now exist in the
    literature, but they are rarely compared \emph{together} on the same RAG system with a
    common protocol for accuracy, cost, and reproducibility---leaving practitioners without
    clear adoption guidance.

    \item \textbf{Critical challenges.}
    \begin{itemize}
      \item \emph{Challenge 1:} \textbf{Claim$\leftrightarrow$source alignment.}
      Faithfulness is a property of individual claims; we must segment answers into
      checkable claims and bind each to the cited evidence before any published scorer
      applies.
      \item \emph{Challenge 2:} \textbf{Fair comparison across method families.}
      LLM-judge, NLI, sampling-consistency, and embedding-similarity scorers differ in
      cost, determinism, and granularity; comparing them fairly on one system and protocol
      is non-trivial.
      \item \emph{Challenge 3:} \textbf{Evaluating against non-circular ground truth.}
      We rely on human span-level labels (RAGTruth~\citep{niu2024ragtruth}) rather than a
      second LLM's opinion.
    \end{itemize}

    \item \textbf{Why is the problem open, and prior work.}
    No single published technique dominates: a 2025 review notes LLM-judge scores are the
    most human-correlated yet costly and opaque~\citep{malin2025faithfulnessreview};
    reliability studies show single-shot LLM judgments are
    inconsistent~\citep{schroeder2024trustjudge, yamauchi2026empiricaljudge}; and
    FaithBench~\citep{bao2025faithbench} is built from cases where leading detectors
    \emph{disagree}, with top detectors near $50\%$ accuracy. On the tooling side,
    commercial LLM-application observability platforms such as LangSmith bundle faithfulness
    checking as paid, hosted LLM-judge-style evaluation, which sharpens---rather than
    answers---the practitioner's question of which technique to adopt and at what cost.
    Grounding foundations: RAG~\citep{lewis2020rag} and REALM~\citep{guu2020realm}.

    \item \textbf{State of the art (our view).} For practical faithfulness scoring,
    LLM-as-judge (RAGAS~\citep{es2024ragas}) is the de-facto standard; NLI-based scoring
    (TRUE~\citep{honovich2022true}) is the strongest non-LLM-judge line. Our study
    benchmarks these published options rather than proposing a new one.
\end{itemize}
